\documentclass[lettersize,journal]{IEEEtran}
\usepackage{amsmath,amsfonts,amssymb}
\usepackage{algorithmic}
\usepackage{algorithm}
\usepackage{array}
\usepackage[caption=false,font=normalsize,labelfont=sf,textfont=sf]{subfig}
\usepackage{textcomp}
\usepackage{stfloats}
\usepackage{url}
\usepackage{verbatim}
\usepackage{graphicx}
\usepackage{color}
% \usepackage{appendix}
\usepackage{cite}
\newcommand{\parab}[1]{\noindent\textbf{#1} }

\hyphenation{op-tical net-works semi-conduc-tor IEEE-Xplore} 

\begin{document}

\appendices
% Insert within the manuscript's existing appendix region.
\section{Full Preprocessing-Aware Security Experiment}
\label{app:pa-security-model}

This appendix specifies the experiment summarized in the main text.
All service, oracle and state rules below belong to that same experiment.

\subsection{Ideal Public Binding Service}
The ideal service $\mathcal B$ maintains the binding table
$\mathsf{UseCtx}$ and provides the read-only interface
$\mathsf{ReadBind}(\mathsf{vk},\mathsf{sid})$, which atomically reads
$\mathsf{UseCtx}[\mathsf{vk},\mathsf{sid}]$ and returns
\[
\begin{cases}
\mathsf{Unbound}, & \mathsf{UseCtx}[\mathsf{sid}]=\bot,\\
\mathsf{Bound}(\mathsf{ctx}_0),
 & \mathsf{UseCtx}[\mathsf{sid}]=\mathsf{ctx}_0.
\end{cases}
\]
Reads are linearizable with token-generation, response and burn
transactions and leave the service state unchanged. The first successful
honest response transaction records the full, canonically encoded context
before releasing the response. This binding remains unchanged after
token-state erasure or identifier burning. Records are indexed by
$(\mathsf{vk},\mathsf{sid})$ and retained for as long as signatures under
$\mathsf{vk}$ may be verified. Only the prescribed state transitions can
update these records; this restriction also applies to corrupted signers.

The public read interface returns the full context of a bound identifier
without exposing private token state or secret shares. The verifier
requires an authenticated, unambiguous read satisfying the service's
consistency requirements. If it cannot obtain such a read, it returns $0$
rather than treating the identifier as $\mathsf{Unbound}$.

Since a binding cannot change, previously authenticated bound records may
be cached. An $\mathsf{Unbound}$ result requires a fresh read, since the
identifier may have been bound since an earlier query. Replacing full
contexts with digests requires an additional collision-resistance argument
to extend the theorem to that implementation.

Let $\mathsf{Vf}_{\mathsf{lat}}$ denote only the canonical-decoding,
challenge and norm checks of the original lattice verifier. When these
checks pass, let $\mathsf{ctx}^{\star}$ be its reconstructed context and
$b_{\mathsf{read}}$ the result of one binding read. The state-assisted
verifier accepts exactly when
\[
\begin{gathered}
\mathsf{Vf}_{\mathsf{lat}}(\mathsf{vk},\mathsf{msg},\sigma)=1,\\
b_{\mathsf{read}}=\mathsf{Unbound}
\quad\text{or}\quad
b_{\mathsf{read}}=\mathsf{Bound}(\mathsf{ctx}^{\star}).
\end{gathered}
\]
The binding check compares the signature context with the service state at
the read's linearization point. Since this permits both unbound
identifiers and matching contexts with incomplete response collections,
unforgeability in these cases still relies on the lattice assumptions.

\subsection{Adversary, Oracles and Shared State}
The model uses two independent random oracles:
\[
\begin{aligned}
\mathcal O_{H_{\mathsf{label}}}&:\{0,1\}^{*}\to R_q,\\
\mathcal O_{H_{\mathsf{ch}}}&:\{0,1\}^{*}\to\mathcal C_{\mathsf{ch}}.
\end{aligned}
\]
The challenger implements each oracle by lazy sampling. A first query
obtains an independent uniform output in the corresponding range, and a
repeated query receives the recorded output. The experiments expose
ordinary random oracles with classical access.

The adversary is a two-stage PPT algorithm
$\mathcal A=(\mathcal A_0,\mathcal A_1)$. Before key generation, it runs
\[
(\mathcal C,\mathsf{st}_{\mathcal A})\leftarrow
\mathcal A_0(1^\lambda,t,n,a_{\max}).
\]
The experiment returns $0$ unless $\mathcal C\subseteq[n]$ and
$|\mathcal C|<t$. Set $\mathcal H=[n]\setminus\mathcal C$. The challenger
then runs
\[
(\mathsf{vk},\{\mathsf{sk}_j\}_{j\in[n]})\leftarrow
\mathsf{KGen}(1^\lambda,t,n,a_{\max}).
\]
$\mathcal A_1$ receives $\mathsf{st}_{\mathcal A}$, $\mathsf{vk}$ and
$\{\mathsf{sk}_j\}_{j\in\mathcal C}$ and may adaptively interleave the
two random-oracle interfaces with the following oracles:

\begin{enumerate}
\item $\mathcal O_{\mathsf{Token}}(j,\mathsf{sid})$ creates, records and
returns the public token of an honest signer. The first successful query
for an unseen identifier initializes it as fresh; later token generation
is permitted only while it remains fresh, and a repeated query for the
same pair returns the recorded token.

\item $\mathcal O_{\mathsf{Resp}}
(j,\mathsf{sid},\mathsf{act},\mathsf{Tok}_{\mathsf{act}},\mathsf{msg})$
first checks structural validity and $j\in\mathcal H\cap\mathsf{act}$,
then locates the shared record and an unconsumed state
$\mathsf{st}_{j,\mathsf{sid}}$ under a fresh or activated identifier.
It checks that the entry indexed by $j$ equals its recorded token,
computes $C_{\mathsf{act}}$ and $c$, and requires the resulting context
to equal the recorded context if already activated. On success, it
atomically binds a fresh identifier, consumes the signer's private state
and releases at most one newly generated honest response for
$(j,\mathsf{sid})$.

\item $\mathcal O_{\mathsf{Bind}}(\mathsf{sid})$ returns
$\mathsf{ReadBind}(\mathsf{vk},\mathsf{sid})$.

\item $\mathcal O_{\mathsf{Burn}}(\mathsf{sid})$ burns the identifier and
erases all remaining honest private states associated with it.
\end{enumerate}

The adversary can locally generate the tokens and responses of corrupted
signers using their secret keys. Structural validity checks concern
canonical encoding, types, ranges, signer membership, and the domain and
entry format of the token mapping. They do not include equality with a
recorded token or consistency with an aggregate, challenge or bound
context. A structurally invalid request returns $\bot$ without changing
state. A response request under an unseen or burned identifier, or for
which no live, unconsumed honest signer state exists, likewise returns
$\bot$ without changing state. Once such a state has been located under
a fresh or activated identifier, a failed recorded-token or
context-consistency check, or any response-generation failure,
atomically burns the identifier and erases all remaining honest private
state before returning $\bot$. These rules also govern aggregate and
challenge consistency checks in $\mathsf{ShareSign}$. Token queries are
idempotent while the identifier is fresh and are rejected after activation
or burning. The burned state is terminal, while previously published
tokens remain public.

The challenger maintains
\[
\mathsf{SidState}[\mathsf{sid}]\in
\{\mathsf{unseen},\mathsf{fresh},\mathsf{activated},\mathsf{burned}\},
\]
$\mathsf{UseCtx}[\mathsf{sid}]\in
\{\bot\}\cup\{(\mathsf{sid},\mathsf{act},C_{\mathsf{act}},\mathsf{msg})\}$,
and $\mathsf{RespSet}[\mathsf{ctx}]\subseteq\mathcal H$.
The experiment fixes one verification key, so the table notation
suppresses it. In a multi-key implementation, every identifier table and
private-state record is indexed by $(\mathsf{vk},\mathsf{sid})$.
Initially, $\mathsf{SidState}[\mathsf{sid}]=\mathsf{unseen}$ and
$\mathsf{UseCtx}[\mathsf{sid}]=\bot$ for every identifier, and
$\mathsf{RespSet}[\mathsf{ctx}]=\varnothing$ for every context.

The first successful honest response request atomically changes a fresh
identifier to activated and records its complete context. Further honest
signers may respond under the same identifier only for that recorded
context, and each pair $(j,\mathsf{sid})$ produces at most one response.
A first successful response also erases every honest private state under
that identifier whose signer is outside the selected active set. A
recorded-token or context-consistency failure after locating a live state,
a response-generation failure, or a burn request moves the identifier to
burned and erases every unconsumed honest private state. Producing one
honest response does not by itself burn the identifier. After a successful
response, the challenger inserts $j$ into
$\mathsf{RespSet}[\mathsf{UseCtx}[\mathsf{sid}]]$.

These tables store the state shared by the token and response oracles and
by honest executions of $\mathsf{TokenGen}$ and $\mathsf{ShareSign}$.
Token creation and retrieval are atomic. To issue a response,
$\mathsf{ShareSign}$ checks the state, establishes or checks the context
binding, consumes the signer's one-time state and releases the response
in a single atomic transaction.

\subsection{Terminal Convention, Authorization and Completion}
Eventually, $\mathcal A_1$ returns one pair
$(\mathsf{msg}^{\star},\sigma^{\star})$. Both winning events are evaluated
on this same terminal output and in the same execution. If
$\mathsf{Vf}^{\mathcal B}(\mathsf{vk},\mathsf{msg}^{\star},
\sigma^{\star})=1$, let
$(\mathsf{sid}^{\star},\mathsf{act}^{\star},C^{\star})$ denote the
identifier, active set and aggregate value reconstructed internally by
verification from the canonical signature encoding. If canonical decoding
fails or verification rejects, these values are undefined and neither
winning event can occur. All probabilities are taken over key and token
generation, the two random oracles, honest protocol executions and
$\mathcal A$.

The model assumes static corruption, reliable erasure and a trusted,
linearizable binding service with authenticated public reads, immutable
bindings and rollback protection. The experiments exclude adaptive
corruption, exposure of honest signers' private states and corruption of
the service. Once $\mathcal A$ produces its terminal output, the protocol
state is frozen, and the winning predicates and final verification are
evaluated against the same values of $\mathsf{UseCtx}$ and
$\mathsf{RespSet}$.

For $\mathsf{ctx}=(\mathsf{sid},\mathsf{act},C_{\mathsf{act}},\mathsf{msg})$,
define
\[
\mathsf{Authorized}(\mathsf{ctx})=1
\iff\mathcal H\cap\mathsf{act}\subseteq\mathsf{RespSet}[\mathsf{ctx}].
\]
Because $\mathsf{act}$ is the exact contributing set, the remaining
shares can be generated locally by the corrupted signers. Under the same
final-state convention, define
\[
\begin{aligned}
\mathsf{Complete}(\mathsf{sid})=1
\iff{}&\exists\mathsf{ctx}:\ \mathsf{sid}(\mathsf{ctx})=\mathsf{sid},\\
&\mathsf{Authorized}(\mathsf{ctx})=1.
\end{aligned}
\]
The quantifier ranges only over recorded contexts. Completion means that
all honest responses required by such a context have been released.
Once true, it remains true after state erasure or identifier burning.

\subsection{Security Notions}
\paragraph{Adaptive-message security}
The adversary may obtain public tokens before choosing its target message.
It wins if it outputs $(\mathsf{msg}^{\star},\sigma^{\star})$ such that
$\mathsf{Vf}^{\mathcal B}(\mathsf{vk},\mathsf{msg}^{\star},
\sigma^{\star})=1$ and there is no recorded context $\mathsf{ctx}$ with
\[
\begin{gathered}
\mathsf{Authorized}(\mathsf{ctx})=1,\\
\mathsf{msg}(\mathsf{ctx})=\mathsf{msg}^{\star}.
\end{gathered}
\]
Partial-response queries make a message nonfresh only when they complete
the honest-response requirements of some recorded context for that message.

\paragraph{One-time-token security}
Suppose verification accepts $(\mathsf{msg}^{\star},\sigma^{\star})$ and
internally reconstructs
$(\mathsf{sid}^{\star},\mathsf{act}^{\star},C^{\star})$. Define
$\mathsf{ctx}^{\star}=(\mathsf{sid}^{\star},\mathsf{act}^{\star},
C^{\star},\mathsf{msg}^{\star})$ and
$\mathsf{ctx}_0=\mathsf{UseCtx}[\mathsf{sid}^{\star}]$.
The one-time-token winning event occurs if
\[
\begin{gathered}
\mathsf{ctx}_0\ne\bot,\qquad
\mathsf{ctx}_0\ne\mathsf{ctx}^{\star},\\
\mathsf{RespSet}[\mathsf{ctx}_0]\ne\varnothing,\\
\mathsf{Vf}^{\mathcal B}(\mathsf{vk},\mathsf{msg}^{\star},
\sigma^{\star})=1.
\end{gathered}
\]
If canonical decoding fails or verification rejects, the event does not
occur. Protection begins as soon as one honest response exposes
information derived from the private state; a complete authorized
transcript for $\mathsf{ctx}_0$ is not required.

\paragraph{Preprocessing-aware unforgeability}
The adversary may arbitrarily interleave all permitted oracle queries.
Let $\mathsf{Win}_{\mathsf{AM}}$ and $\mathsf{Win}_{\mathsf{OT}}$ denote
the adaptive-message forgery and incompatible identifier-reuse events,
respectively. It wins the preprocessing-aware experiment if
\[
\mathsf{Win}_{\mathsf{PA}}=
\mathsf{Win}_{\mathsf{AM}}\lor\mathsf{Win}_{\mathsf{OT}}.
\]
These events use $\mathsf{Vf}^{\mathcal B}$, with the PA advantage also
denoted by
$\operatorname{Adv}^{\mathsf{PA},\mathcal B}_{\Pi_{\mathsf{pp}},\mathcal A}$.
The binding check provided by $\mathcal B$ rules out the
incompatible-context event.

For $X\in\{\mathsf{AM},\mathsf{OT},\mathsf{PA}\}$, define
\[
\operatorname{Adv}^{X}_{\Pi_{\mathsf{pp}},\mathcal A}(\lambda)
=\Pr[\mathsf{Win}_X].
\]
The corresponding notion holds if this advantage is negligible for every
PPT adversary satisfying the corruption bound. The AM branch is ordinary,
non-strong, message-based unforgeability: once any context for a message
is fully authorized, a second accepting signature on that message is not
an AM win. The OT branch separately excludes accepting reuse of an
activated identifier under an incompatible context. Both branches are
evaluated on the same terminal signature.

\end{document}
